\documentclass[a4paper, 11pt]{article}

% --- Margini di pagina (meno larghi rispetto alla tesi) ---
\usepackage[a4paper, margin=2.5cm]{geometry}

% --- Pacchetti tipografici e lingua ---
\usepackage[utf8]{inputenc}
\usepackage[T1]{fontenc}
\usepackage{lmodern}
\usepackage{newunicodechar}
\newunicodechar{’}{'}
\newunicodechar{Ł}{\L{}}
\usepackage{setspace}
\usepackage{microtype}

% --- Matematica e simboli ---
\usepackage{amsmath}
\usepackage{amssymb}

% --- Figure e tabelle ---
\usepackage{graphicx}
\graphicspath{ {./images/} }
\usepackage{booktabs}
\usepackage{tabularx}
\usepackage{adjustbox}
\usepackage{subcaption}
\usepackage{float}
\usepackage{enumitem}

% --- Collegamenti ipertestuali (stile come da modello allegato con
% box verde per le citazioni) ---
\usepackage{xcolor}
\usepackage[
  colorlinks=false,
  citebordercolor={0 1 0},
  linkbordercolor={0.8 0 0},
  urlbordercolor={0 0.5 1},
  pdfborder={0 0 1}
]{hyperref}

% --- Bibliografia condivisa con la tesi (biblatex + biber, stile IEEE) ---
\usepackage[
  backend=biber,
  style=ieee,
  sorting=none,
]{biblatex}
\addbibresource{bibliografia.bib}

% --- Soppressione numerazione sezioni per documento di sintesi ---
\setcounter{secnumdepth}{0}

\pagestyle{plain}

\begin{document}

% ==============================================================================
% INTESTAZIONE PRIMA PAGINA (Titolo e Partecipanti come da modello)
% ==============================================================================
\begin{center}
  {\LARGE \textbf{A Language-Agnostic Framework\\ for Dependency
  Extraction}\par}

  \vspace{1.6em}
  {\large \textbf{Candidate:} Christian Ferrario\par}

  \vspace{1.0em}
  {\large
    \textbf{Supervisor:} Prof. Francesca Arcelli Fontana\\[0.3em]
  \textbf{Co-supervisor:} Dr. Francesco Refolli\par}

  \vspace{1.6em}
  {\large October 2026\par}
\end{center}

% ==============================================================================
% 1. OBJECTIVES OF THE STUDY
% ==============================================================================
\section{Introduction and Objectives of the Study}
\label{sec:objectives}

Modern software systems undergo continuous evolution to adapt to
changing requirements, often leading to architectural degradation and
accumulated \textbf{technical debt} \cite{lehman_programs_1980,
kruchten_technical_2012}.
While low-level code smells primarily degrade local code quality
\cite{fowler_refactoring_2013}, architectural flaws such as cyclic
dependencies and high coupling (\textit{Architectural Smells})
severely hinder system maintainability and impede safe refactoring operations
\cite{garcia_identifying_2009, arcelli_fontana_independent_2019}.

To detect and prevent architectural degradation, software organizations
rely on automated static analyzers that construct \textbf{Dependency
Graphs}, modeling components as nodes and their semantic relationships
as directed edges.
However, existing analyzers suffer from a fundamental limitation:
\textbf{they are tightly coupled to specific programming languages or
compiler front-ends}.
Supporting multiple languages typically requires integrating
dedicated compiler front-ends and re-engineering the extraction
pipeline for each language, incurring substantial engineering overhead.
Cross-language inspection of Concrete Syntax Trees (CSTs) generated by
parsers such as Tree-sitter \cite{brunsfeld_treesitter_2018} reveals
remarkable recurring structural similarities across divergent languages.
This thesis exploits these structural commonalities to design and
implement an architectural dependency extraction pipeline that is
inherently language-agnostic.

\subsection{Research Objectives and Scope}
The primary objective of this thesis is to design, formalize, and
implement \textbf{OmniDeps}, a language-agnostic static analysis
framework capable of extracting comprehensive architectural dependency
graphs across heterogeneous programming languages (specifically targeting
C, C++, Java, Python, and Rust) through a single unified pipeline.

Specifically, the research pursues three core objectives:
\begin{itemize}[leftmargin=2.0em, noitemsep]
  \item \textbf{Design and Formalization}: Formalizing an
    Intermediate Representation (IR) through strongly typed algebraic data
    types \cite{pierce_types_2002} and a query model for name
    resolution, abstracting architectural declarations and behavioral
    dependencies away from language-specific syntax.
  \item \textbf{Framework Implementation}: Engineering a
    deterministic pipeline that transforms native Tree-sitter CSTs
    into normalized dependency graphs through heuristic syntactic
    extraction, query-based name resolution, and architectural graph projection.
  \item \textbf{Empirical Validation}: Verifying extraction
    capabilities against ground-truth semantic oracles across
    language constructs, and benchmarking cross-tool structural
    alignment and scalability against established analyzers
    (\textit{Arcan}~\cite{fontana_arcan_2017},
      \textit{Depends}~\cite{depends_github}, and
    \textit{Entangle}~\cite{teruzzi_2024}) on real-world repositories
    using a formal graph harmonization methodology.
\end{itemize}

To achieve this without re-implementing an entire compiler toolchain,
the framework operates under the foundational premise that \textbf{the
  analyzed source code compiles successfully within its native
toolchain}.
This deliberately scopes out compiler-level validity checks (such as
type checking or visibility enforcement) and enables OmniDeps to focus
strictly on architectural dependency extraction.

\subsection{Methodology}
To guarantee extensibility, formal soundness, and cross-language
consistency, OmniDeps structures the analysis into three sequential,
deterministic phases.

\paragraph{Syntactic Extraction to Universal IR}
Rather than maintaining dedicated parser front-ends for each target
language, the extraction stage exploits the structural convergence
observed in Tree-sitter CSTs across different programming languages.
Lightweight heuristic classification predicates inspect syntax tree
nodes to extract architectural entities directly into our strongly typed
universal IR, formalized through algebraic data types \cite{pierce_types_2002}.
This stage captures:
\textit{modular boundaries} (packages, namespaces, modules);
\textit{structured types} (classes, structs, interfaces, traits, and
enums) along with their inheritance links and generic type parameters;
\textit{behavioral members} (functions, methods, constructors, and
fields); and \textit{auxiliary declarations} (such as explicit and wildcard
imports, type aliases, decorators).
Statement blocks inside function bodies are also extracted to
capture intra-procedural operations, including method calls, type
instantiations, and local variable declarations.
By mapping heterogeneous concrete syntax nodes into algebraic sum and
product types, the framework abstracts away language-specific keywords,
grammars, and syntax rules at the very entrance of the pipeline,
allowing all subsequent phases to operate exclusively on standardized
abstractions.

\paragraph{Name Resolution via Query Algebra}
Resolving symbolic identifiers across heterogeneous languages presents a
critical architectural challenge: within function and method bodies,
references depend heavily on transient procedural context (such as
  local variables, formal parameters, lexical shadowing, and implicit
receivers like \texttt{this} or \texttt{self}).
Attempting to resolve all dependencies in a single global step would
couple local intra-procedural state across every block in the workspace
with global resolution, leading to state pollution and
order-dependent complexities.
To resolve this issue, OmniDeps decomposes name resolution into two
decoupled, sequential sub-phases:
first, a \textit{Query Building} pass traverses method bodies using an
ephemeral Symbol Stack to eliminate local variables and parameters,
rewriting references into an algebraic query language governed by three
orthogonal primitives: \texttt{Find} (ascending lexical scope climbing),
\texttt{Extract} (structural member descent across inheritance
hierarchies), and \texttt{Call} (return-type propagation in chained
invocations).
Second, a \textit{Query Execution} pass evaluates these self-contained
queries deterministically against a global \textbf{Scope Tree} capturing
the modular containment of the entire workspace.
Cross-file references, imports, and inheritance hierarchies are resolved
independently of source file processing order, binding symbols to
internal definitions, external dependencies, or language primitives.

\paragraph{Dependency Graph Construction}
The final phase projects the resolved IR into an architectural directed
multigraph.
First, the hierarchical module tree is recursively flattened to promote
software components to independent graph vertices, each assigned a
qualified name computed from its ancestor module paths.
Then, a traversal pass inspects all resolved type references to
instantiate directed dependency edges across three core categories:
\textit{structural containment} (module, component
nesting);
\textit{type-level couplings} (field declarations, method parameter and
return types, type aliases, and generic bounds); and
\textit{behavioral interactions} (method invocations, object
instantiations, and attribute accesses extracted from function bodies).
To preserve referential integrity, references targeting verified external
libraries or built-in types dynamically synthesize corresponding
\texttt{External} and \texttt{Primitive} endpoints, while unresolvable
references are discarded.
Finally, duplicate edges between identical endpoints are pruned to
prevent distorted coupling metrics, producing a clean, normalized
dependency graph ready for architectural evaluation.

% ==============================================================================
% 2. EVALUATION AND RESULTS
% ==============================================================================
\section{Evaluation and Results}
\label{sec:results}

The empirical validation of OmniDeps is structured around four research
questions covering extraction correctness across language constructs,
cross-tool structural alignment on real-world repositories, discrepancy
decomposition across dependency types, and computational scalability.

\subsection{Synthetic Ground-Truth Benchmark Validation}
To evaluate extraction correctness in a controlled setting, OmniDeps
was validated against manually verified architectural
oracles across all five supported languages (C, C++, Java, Python,
and Rust). These benchmarks served during development as a regression
test suite to exercise individual language constructs and verify the
baseline behavior of the extraction pipeline.

OmniDeps passes all benchmark test cases, recovering the expected
entities and dependencies across all five target languages. However,
synthetic micro-benchmarks are limited in scope and cannot guarantee
generalization to arbitrary, large-scale software projects, which
motivates the empirical evaluation on real-world repositories.

\subsection{Comparison with State-of-the-Art Analyzers}
Comparing results across independent tools is
challenging due to divergent graph representations and granularity
levels. To enable empirical comparison, we formulated a graph
harmonization framework that normalizes entity identifiers through
qualified paths, lifts fine-grained members to enclosing
architectural types, synthesizes package containment hierarchies, and
maps tool-specific edge labels onto a unified taxonomy.

Using this framework, OmniDeps was evaluated against three tools
(Depends, Arcan, and Entangle) across a dataset of \textbf{156
repositories} comprising nearly 98 million lines of code.

\paragraph{Analysis Completion.}
OmniDeps achieved an overall completion rate of \textbf{99.4\%}
(155/156 projects successfully analyzed). Its sole failure occurred
on MongoDB due to a stack overflow during AST traversal. In contrast,
baseline tools exhibited lower completion rates under a generous
one-hour timeout ($T_{\max} = 3600\,\text{s}$): Depends completed
92.7\% of projects (115/124), failing primarily due to timeouts on
large codebases and on a syntax parsing error on Checkstyle; Arcan
reached 49.2\% (61/124), exiting prematurely due to execution errors;
and Entangle completed 29.7\% (19/64), failing on Rust due to
unhandled runtime exceptions.

\begin{table}[htpb]
  \centering
  \small
  \caption{Pairwise structural similarity (mean Smoothed Jaccard
  Index) across real-world repositories.}
  \label{tab:results_summary}
  \vspace{0.4em}
  \begin{adjustbox}{max width=\textwidth}
    \begin{tabular}{lcccccc}
      \toprule
      \textbf{Language} & \multicolumn{2}{c}{\textbf{OmniDeps vs.
      Depends}} & \multicolumn{2}{c}{\textbf{OmniDeps vs. Arcan}} &
      \multicolumn{2}{c}{\textbf{OmniDeps vs. Entangle}} \\
      \cmidrule(lr){2-3} \cmidrule(lr){4-5} \cmidrule(lr){6-7}
      & Nodes & Edges & Nodes & Edges & Nodes & Edges \\
      \midrule
      \textbf{Java}             & 0.814 & 0.587 & 0.630 & 0.580 &
      0.770 & 0.600 \\
      \textbf{Python}           & 0.740 & 0.296 & ---   & ---   & ---
      & ---   \\
      \textbf{C (Raw)}          & 0.191 & 0.115 & ---   & ---   & ---
      & ---   \\
      \textbf{C (Preprocessed)} & 0.204 & 0.136 & ---   & ---   & ---
      & ---   \\
      \textbf{C++ (Raw)}        & 0.143 & 0.081 & ---   & ---   & ---
      & ---   \\
      \textbf{C++ (Preprocessed)} & 0.192 & 0.165 & --- & ---   & ---
      & ---   \\
      \textbf{Rust}             & ---   & ---   & ---   & ---   & ---
      & ---   \\
      \bottomrule
    \end{tabular}
  \end{adjustbox}
\end{table}

As reported in Table \ref{tab:results_summary}, in Java OmniDeps
aligns closely with all three baseline tools, with edge Jaccard
similarity around 0.59. OmniDeps exhibits low
entity omission rates: 2.1\% against Depends
($\text{Inf}_{\text{nodes}} = 0.021$) and 0.3\% against Arcan
($\text{Inf}_{\text{nodes}} = 0.003$).
Decomposing edge alignment by relationship type reveals that containment
hierarchies (\texttt{belongsTo}) exhibit near-perfect agreement
($\ge 98.2\%$, omission $\le 1.8\%$) and inheritance links match closely.
Discrepancies concentrate predominantly in behavioral invocations
(\texttt{dependsOn}): while OmniDeps resolves return types along local
call chains, it does not load the standard library and
external binaries, which compiler-backed tools use to trace chained calls
traversing external types.

In Python, node similarity with Depends is high (0.740). While OmniDeps
leverages \textit{type annotations} to resolve call targets,
real-world codebases adopt type hints only partially
or inconsistently \cite{magalhaes_understanding_2026}, leading to
lower edge agreement (0.296). For Rust, baseline tools provide no
architectural support, leaving
pairwise comparison unavailable; nevertheless, OmniDeps successfully
analyzed all 32 Rust repositories at over 97,000
LOC/s.

\paragraph{Structural Alignment in C and C++: Arcan and Depends.}
Comparing OmniDeps against existing tools in C and C++ highlights two distinct
phenomena. First, comparison with Arcan results in zero overlap due to
divergent abstractions: Arcan models procedural code at the file and
header level for smell detection, whereas OmniDeps and Depends
extract structured types, producing disjoint entity sets.
Second, the moderate Jaccard similarity observed with Depends
(mean 0.14 to 0.20) stems from preprocessor mechanics and scope differences
rather than extraction deficiencies. In raw source code, extensive macros mask
type declarations from syntax parsers. Under translation-unit preprocessing,
macro expansion improves extraction, but inlining introduces standard library
headers that OmniDeps extracts while Depends excludes. Furthermore, OmniDeps
models enum variants as individual architectural entities while Depends prunes
thems. These factors inflate node counts, mathematically
suppressing the Jaccard denominator. Crucially, entity omission remains low
($\text{Inf}_{\text{nodes}} \le 0.202$), confirming that OmniDeps captures
virtually all core project entities. On cleanly modular codebases without
massive header inlining, concordance is high: OmniDeps achieves identical graphs
on \texttt{sds} ($J = 1.0$) and strong edge agreement on systems such as
\texttt{ninja} ($0.410$) and \texttt{re2} ($0.459$).

\subsection{Performance and Scalability}
Across the evaluation dataset, OmniDeps exhibits execution runtimes
that scale essentially linearly with codebase size. Analysis reliably
completes within the order of seconds for the vast majority of
repositories, scaling predictably without runtime spikes even on
systems exceeding hundreds of thousands of lines of code.
When compared to established static analyzers, OmniDeps achieves
execution times that are consistently comparable to or faster than
baseline tools. Sustaining an average throughput of 74,445 LOC/s, OmniDeps
provides average speedups of \textbf{13.2x over Depends}, \textbf{48.2x over
Arcan}, and \textbf{4.0x over Entangle}.
These results confirm that lightweight syntax-tree extraction avoids
the significant runtime overhead of compiler front-ends while
maintaining practical scalability for continuous integration pipelines.

\subsection{Threats to Validity}
In C and C++, translation-unit preprocessing inlines standard library headers,
which together with fine-grained enum variant extraction and vendored subtrees
inflates node counts and suppresses Jaccard similarity against compiler-backed
tools. In addition, the pipeline currently relies on manually curated catalogs
to identify built-in primitive types, and does not model C-style variadic
functions with untyped ellipsis parameters. Finally, the empirical validation
covers imperative and object-oriented languages within the C family, leaving
the applicability of the intermediate representation and name resolution
pipeline to functional or declarative paradigms unverified.

\subsection{Conclusions and Future Work}
This thesis demonstrates that language-agnostic architectural
dependency extraction is feasible, sound, and scalable, unifying
multi-language parsing through a common intermediate representation
and scoped queries.
As the thesis focused primarily on model design and extraction
strategies rather than low-level tuning, future work includes
introducing multithreading to parallelize extraction. In addition,
while current support centers on procedural and object-oriented
C-like languages, the framework could be extended to functional or
declarative paradigms, as well as to detecting cross-language
dependencies like JNI or FFI calls across polyglot systems.

% ==============================================================================
% BIBLIOGRAFIA CONDIVISA
% ==============================================================================
\printbibliography[title={References}]

\end{document}
